\section{人形机器人}

\subsection{双足行走与全身控制}

讲者简要介绍了 RL 在人形机器人上的应用进展：
\begin{itemize}
    \item 双足稳定行走
    \item 全身协调（同时走路和操作）
    \item 对外力扰动的鲁棒性
\end{itemize}

\begin{knowledgebox}{人形机器人的特殊挑战}
相比四足机器人，人形机器人的 RL 训练面临更大挑战：
\begin{itemize}
    \item 不稳定的双足支撑（更窄的稳定裕度）
    \item 更高的自由度（全身 30+ 关节）
    \item 上肢和下肢需要协调
    \item 跌倒的后果更严重（硬件昂贵）
\end{itemize}
\end{knowledgebox}

\subsection{本章小结}

人形机器人代表了 RL for robotics 的前沿挑战，当前进展快速但离稳定部署仍有距离。

